\noindent\textbf{Datasets.} Both benchmark datasets originate from the C-LLM evaluation~\cite{xian2024cllm} and were automatically translated from Chinese to English. \emph{MIX} contains 100~open-ended questions spanning ten domains (science, history, mathematics, language, culture, geography, entertainment, technology and personal/conversational topics). \emph{PRO} contains 60~four-option multiple-choice physics questions stratified across three educational levels: Junior High School (20), High School (20) and University (20). Both are stored as JSON arrays with fields \texttt{question} and \texttt{answer}. Binary yes/no variants used by the DeepThought baseline are in \texttt{deepthought-comparison/datasets/}.

\smallskip\noindent\textbf{Experimental Configurations.} Table~\ref{tab:configs} summarizes the 64~configurations evaluated in the paper. Each configuration is uniquely identified by the combination of LLM model, decoding setting, dataset and node split. The evaluated LLMs are ChatGPT-4o-mini, Gemini-2.5-flash-lite (configured with 8192 output tokens and reasoning disabled) and DeepSeek-v4-flash in two inference modes: \emph{non-thinking} (referred to as chat) and \emph{thinking} (referred to as reasoner). The directory naming convention maps directly to these parameters.

\begin{table}[t]
\centering
\caption{Experimental configurations (64 total = 4 $\times$ 4 $\times$ 2 $\times$ 2).}
\label{tab:configs}
\footnotesize
\begin{tabular}{llcl}
\toprule
\textbf{Dimension} & \quad\textbf{Values} & \textbf{Count} \\
\midrule
LLM family & \quad GPT-4o-mini, Gemini-2.5-flash-lite, & 4 \\
           & \quad DeepSeek-v4-flash (chat, reasoner) & \\
Decoding   & \quad $C_1$: $T{=}0$; $C_2$: $T{=}0.5$; & 4 \\
           & \quad $C_3$: $T{=}1.0$; $C_4$: $T{=}1.0$, seed$=4321$ & \\
Dataset    & \quad MIX (100\,Qs, 30 runs), PRO (60\,Qs, 20 runs) & 2 \\
Node split & \quad 6/4 (60-40), 7/3 (70-30), $N{=}10$ & 2 \\
\bottomrule
\end{tabular}
\end{table}


\smallskip\noindent\textbf{LLM Responses.} For each configuration, the repository includes the complete set of LLM-generated answers in JSON format. Each entry in \texttt{q\_*\_answers.json} contains a question string and an array of 10~answers: the first $N$~are independently generated honest responses, the remaining $10{-}N$ are identical copies of the coordinated malicious response produced via the adversarial prompt (Figure~2 in the original paper).

\smallskip\noindent\textbf{Credibility Logs.} The files \texttt{node\_weights\_log\_*.txt} record the evolution of node credibility across questions. Each line contains 10~floating-point values representing the updated credibility of each node after processing one question. The credibility update follows the C-LLM procedure: BERT embeddings (\texttt{bert-base-uncased}, mean pooling) are used to compute pairwise cosine similarity, which is then combined with current node weights to produce updated weights. All nodes are initialized at $0.50$, as per the original C-LLM evaluation~\cite{xian2024cllm} .

\smallskip\noindent\textbf{Analysis Reports.} The \texttt{results/} directory contains 16~Excel workbooks (one per model $\times$ decoding configuration) plus a cross-model dashboard. Each workbook includes sheets for single-run weight evolution, shuffled-run weight evolution, per-run accuracy and credibility delta sequences. Statistical summaries (FCD with 95\% CI, SCD incidence) are in the \texttt{final-deltas/} and \texttt{systemic-dominance/} subdirectories.

\smallskip\noindent\textbf{SBERT Model.} The credibility computation uses \texttt{bert-base-uncased}, consistent with the Sentence-BERT approach described in Reimers and Gurevych~\cite{reimers2019sbert}. Sentence embeddings are obtained via mean pooling over the last hidden state. The model is automatically downloaded (${\sim}$440\,MB) on first execution of Step~3.
